
% =============================================================================
% Section 05 — Invariance and Difference
% Book I — Ontology of Viability
% Chapter 02 — Invariance
%
% Purpose:
%   To examine the relation between invariance and difference,
%   establishing how preservation may coexist with difference
%   within and across transformed configurations.
%
% Conceptual Policy:
%   1. Invariance does not eliminate difference.
%   2. Difference does not necessarily eliminate invariance.
%   3. Preservation and difference must be assessed with respect
%      to specified determinations, relations, and conditions.
%   4. Similarity, equivalence, and identity must not be conflated.
%   5. The persistence of an invariant does not establish
%      the absence of ontological difference.
%   6. Difference must not be reduced to variation in a single
%      identifiable feature.
%
% =============================================================================

\section{Invariance and Difference}
\label{sec:invariance-and-difference}

\subsection{The Ontological Problem of Invariance and Difference}

The relation between invariance and difference constitutes
a further problem in the ontological analysis of preservation.

Invariance concerns the preservation of a specified
determination or relation across a transformation.
Difference concerns the distinction or non-identity
between determinations, configurations, or existents.

These concepts address different aspects of ontological
reality and must not be treated as mutually exclusive.

The preservation of a particular feature does not
necessarily eliminate differences between the configurations
in which that feature obtains. Similarly, the existence
of difference does not necessarily imply that every
determination or relation has been altered.

An existent may exhibit invariant features while
undergoing differences in other respects. Different
existents may also share particular invariant features
without thereby becoming numerically identical.

The central problem is therefore to clarify how preservation
and difference can coexist without reducing either
concept to the other.

\subsection{Invariance Does Not Eliminate Difference}

The preservation of a determination or relation does
not establish the absence of difference.

A particular feature may remain preserved while other
determinations, relations, or configurations undergo
alteration.

Consequently, two configurations may share an invariant
and nevertheless differ in several other respects.

The preservation of a structural organization, for example,
does not necessarily imply the preservation of every
quantitative or qualitative determination associated
with that organization.

Likewise, the preservation of a particular relation
does not establish that the relata participating
in that relation are identical in every respect.

Invariance must therefore be understood as the preservation
of a specified feature rather than as the elimination
of all difference.

To infer the absence of difference from the existence
of an invariant would be to extend a restricted
preservation claim beyond its legitimate scope.

\subsection{Difference Does Not Necessarily Eliminate Invariance}

The converse is equally important.

The existence of difference between two configurations
does not necessarily establish that no invariant
is preserved between them.

A transformation may produce differences in certain
determinations while leaving other features intact.

The configurations before and after the transformation
may therefore be different without being entirely
discontinuous in every relevant respect.

The same possibility applies to relations between
different existents.

Two existents may differ in their determinations,
compositions, or configurations while sharing a
particular relational or structural feature.

Such preservation does not eliminate their differences.
It establishes only that a specified feature remains
common or preserved under the conditions considered.

Difference and invariance must consequently be examined
at the appropriate analytical level.

\subsection{Difference Between Determinations and Difference Between Configurations}

The concept of difference must be considered at more
than one analytical level.

A difference between individual determinations is not
necessarily equivalent to a difference between entire
configurations.

Two configurations may contain different determinations
while preserving a particular organizational relation.

Conversely, two configurations may contain similar
determinations while differing in the way those
determinations are organized or related.

The preservation of individual determinations therefore
does not automatically establish the preservation
of a configuration.

Likewise, the alteration of individual determinations
does not necessarily establish the alteration
of every structural or relational feature.

An adequate account of invariance must distinguish
between the preservation or alteration of individual
determinations and the preservation or alteration
of their configurations.

Without this distinction, different forms of preservation
and difference may be incorrectly conflated.

\subsection{Difference and the Scope of Comparison}

Every assertion of difference presupposes a basis
of comparison.

The relevant basis may concern determinations,
relations, structures, compositions, or other
specified aspects of existents.

Two configurations may be different under one criterion
while remaining equivalent under another.

For example, two structures may differ in their
constituent elements while preserving an equivalent
organizational pattern.

Such equivalence does not establish numerical identity.
It establishes a correspondence under a specified
criterion of comparison.

Similarly, the identification of a difference in
one respect does not establish that the compared
configurations differ in every respect.

The scope of comparison must therefore be specified
before conclusions about difference or preservation
can be established.

This requirement parallels the conditions developed
for identifying invariants, since both concepts
require a determinate subject and respect of analysis.

\subsection{Invariance, Similarity, and Equivalence}

Invariance must also be distinguished from similarity
and equivalence.

Similarity indicates a degree of resemblance between
determinations or configurations.

Equivalence indicates correspondence under a specified
criterion or relation of comparison.

Invariance, by contrast, concerns the preservation
of a specified feature across a transformation
or within a defined domain of comparison.

These concepts may be related, but they are not
interchangeable.

Two configurations may be similar without preserving
the same determination through a transformation.
Two structures may be equivalent under a formal
criterion without being ontologically identical.

Likewise, the preservation of an invariant does
not necessarily imply that the initial and resulting
configurations are similar in every respect.

The distinction is particularly important when
formal relations are used to describe ontological
configurations.

A formal equivalence may establish a correspondence
within a model, but it does not automatically
establish the numerical identity of the existents
represented by that model.

The ontological status of the compared configurations
must therefore remain distinct from the formal
criterion through which their correspondence
is established.

\subsection{Invariance Across Different Existents}

The problem of invariance and difference is not
restricted to transformations occurring within
a single existent.

A particular determination or relation may be
present in different existents.

The presence of a shared feature does not establish
that the existents possessing it are numerically
identical.

Rather, it establishes that a particular feature
or relation is common to them under the relevant
criterion.

The distinction between shared invariance and
numerical identity is essential.

Different existents may exhibit similar structures,
preserve equivalent relations, or possess common
determinations while remaining distinct existents.

The existence of common features therefore cannot
be used as a sufficient basis for reducing
ontological difference to identity.

This also indicates that invariance may be relevant
to comparisons among existents without determining
their numerical identity.

\subsection{Difference Without the Loss of Every Invariant}

The occurrence of difference does not necessarily
eliminate every invariant associated with an existent.

A transformation may alter a particular determination
while preserving a relation or structural feature.

Similarly, the emergence of a new configuration
may involve both the alteration of certain features
and the preservation of others.

The identification of difference must therefore
not be treated as sufficient grounds for denying
all forms of invariance.

However, the converse limitation must also be maintained:
the preservation of one feature cannot establish
that every other feature remains unchanged.

The coexistence of difference and invariance requires
an analysis that remains specific to the relevant
determinations, relations, and conditions.

Neither complete preservation nor complete alteration
should be presumed without sufficient ontological
justification.

\subsection{The Limits of Invariance as a Basis of Distinction}

Although invariant features may contribute to
the differentiation or comparison of existents,
they cannot be treated as an exhaustive basis
for establishing ontological distinction.

Two existents may share particular invariant
features while differing in other determinations
or relations.

Conversely, an existent may undergo alteration
in previously identified invariant features
without thereby losing its identity.

An invariant can therefore provide a basis for
describing a particular aspect of an existent,
but it cannot independently determine the full
ontological status of that existent.

This limitation is consistent with the distinction
between identity and invariance established earlier.

The analysis of difference must consequently
remain open to determinations and relations
that are not captured by a particular invariance
claim.

\subsection{Synthesis}

Invariance and difference are not mutually exclusive
ontological conditions.

The preservation of a specified determination
or relation does not eliminate differences in
other respects. Similarly, the existence of
difference does not necessarily eliminate every
invariant.

Difference may concern individual determinations,
relations, or entire configurations. Its assessment
therefore requires a specified basis and scope
of comparison.

Similarity and equivalence must also remain
distinguishable from invariance and numerical
identity. A shared or preserved feature does
not establish that the existents possessing
it are identical.

Invariance contributes to the analysis of
preservation and comparison, but it cannot
exhaust the ontological problem of difference.

The next section examines the relation between
invariance and structure, addressing whether
the preservation of structural organization
can be distinguished from the preservation
of individual determinations and relations.
